\section{Climate Change Committee scenarios}
\label{appendix:ccc_scenarios}

The Climate Change Committee \citep{cccSixthCarbonBudget2020} provides series of of ``resource cost'' (capex and opex) for each of its 5 scenarios. Here we detail how we translate these into prices and then into productivity. 


\subsection{Framework}

The intuitive idea is that extra capital investment of the net zero transition, net of the savings that it will bring, translates as changes in unit costs, which we may translate one-to-one as changes in unit price. The CCC gives time series for 14 sectors and the 5 scenarios of ``Additional capital investment'', ``Additional operating costs'' and the discounted combination of the two in a ``Annualised resource cost'' time series ($C_{i,t}$ hereafter).\footnote{See the \href{https://www.theccc.org.uk/2021/02/01/the-numbers-behind-the-budget-six-ways-to-explore-the-sixth-carbon-budget-dataset/}{CCC high level piece on the 6th Carbon number} and download the dataset \href{https://www.theccc.org.uk/wp-content/uploads/2021/02/The-Sixth-Carbon-Budget-Dataset.xlsx}{at this link}.}

We make two important assumptions:

First, we assume that OPEX and CAPEX numbers provided by the CCC apply to industries. In practice, the CCC conflates OPEX and CAPEX numbers of firms and households.

Second, we assume that production in volume remains constant. The scenarios do not provide a volume proxy for each sector, and we do not expect large increases in output, so we think this is the right benchmark. We test the influence of this assumption at the end of this section.

The method is as follows. First, we sum the annualised resource cost between $t$ and $t+\tau$ as 
\begin{equation} C_{i,t,\tau}=\sum_{t'=t}^{t+\tau}C_{i,t'}.
\end{equation}
Thus, if the total increase in cost between $t$ and $t+\tau$ is $C_{i,t,\tau}$, the total increase in \textit{unit} cost is $C_{i,t,\tau}/Q_i$. Assuming unit price growth follows unit costs growth, we have
\begin{equation}
    \hat{p}_{i,t,\tau}=\frac{p_{i,t+\tau}-p_{i,t}}{p_{i,t}}=\frac{C_{i,t,\tau}/Q_i}{p_{i,t}}.
    \label{eq:phatfromchat}
\end{equation}
In our theoretical framework\footnote{
%
In our theoretical framework, $\hat{p}$ is a log change, while here we are using a percentage. This introduces a non negligible approximation error for the growth magnitudes considered here. We omit this issue here.
%
}, the productivity growth of the sector $i$ can be predicted as (Eq. \ref{eq:TFP_from_prices_and_output} or
\ref{eq:ai_f_of_zeta_main}),
\begin{equation}
\begin{aligned}
    \hat{A}_{i,t,\tau}&=\zeta \hat{Q}_{i,t,\tau}+(\zeta-1) \hat{p}_{i,t,\tau}\\
    &=(\zeta-1) \frac{C_{i,t,\tau}}{p_{i,t}Q_i},
    \end{aligned}
\end{equation}
where we have used Eq. \ref{eq:phatfromchat} and the assumption that output in volume remains constant, $\hat{Q}_{i,t,\tau}=0$.
using Domar aggregation (Eq. \ref{eq_domar} or \ref{eq:Acttau_pred_main}), aggregate productivity growth then is 
\begin{equation}
\label{eq_agg_prod_CCC}
\begin{aligned}
    \hat{A}_{t,\tau}&=\sum_{i=1}^N \frac{p_{i,t}Q_i}{p_{t}Y_t} \hat{A}_{i,t,\tau}\\
    &=\sum_{i=1}^N \frac{p_{i,t}Q_i}{p_{t}Y_t}(\zeta-1) \frac{C_{i,t,\tau}}{p_{i,t}Q_i}\\
      &= (\zeta-1) \frac{\sum_{i=1}^N C_{i,t,\tau}}{p_{t}Y_t},\\
\end{aligned}
\end{equation}
where $p_{t}Y_t$ is the value of GDP. Note that the quantities level and prices ($p_i$ and $Q_i$) disappear from Eq. \ref{eq_agg_prod_CCC}, and with them disappear the difficulty of defining Domar weights for cross-industries sector, e.g. ``Heat". Thus, the impact on aggregate productivity growth is remarkably simple to compute - it is just the total resource cost as a fraction of today's GDP, multiplied by $\zeta-1 \approx -0.13$. 

%\FL{We need to think more carefully about what happens to the ``non-CCC'' industries. e.g. assume that they collectively represent 25\% of GDP. This should guide our assumption for volume growth}

Below, we compute aggregate productivity growth in CCC scenarios two ways. First, we directly consider the aggregate "Whole economy" resource cost of the transition as given by the CCC scenario. Second, we break down that result into sectoral contributions. The contribution of sector $i$ to aggregate productivity is $(\zeta-1) \frac{C_{i,t,\tau}}{p_{t}Y_t}$.


\subsection{Results}

For Balanced Net Zero, the result is a a negative shock of productivity of 1460bp. All five CCC scenario show productivity decrease between -1970bp and -1070bp (Table \ref{tab_CCC_TFP_opex_capex}).

Figure \ref{fig_cumAhat_CCC_opex_capex} show the dynamics of productivity growth from 2020 and 2050. Table \ref{tab_CCC_TFP_opex_capex} and Figure \ref{df_waterfall_ccc}  break down the contribution of each sector for the 5 CCC scenarios.  Residential and non-residential buildings are driving most of the productivity decrease  (-870bp in Balanced Net Zero). The positive contributions of surface transport, aviation and agriculture do not compensate the residential and manufacturing negative productivity shock. Note that the electricity supply contributes negatively to the productivity growth whereas the cost of electricity should decrease in the future. It is an incentive to build our own scenarios based on the cost and deployment of technologies.


\begin{figure}[H]
    \centering
        \caption{CCC Balanced Net Zero Pathway aggregate productivity growth contribution per sector}
    \label{fig_cumAhat_CCC_opex_capex}
    \includegraphics[width=\textwidth]{figures/figures_CCC/g_cumAhat_CCC_opex_capex.jpeg}
\end{figure}

\input{figures/figures_CCC/tab_CCC_TFP_opex_capex}


\subsection{Robustness}

We have assumed in the previous part that there was no output growth for any of the sectors. We now evaluate how relaxing this assumption changes the results.

The CCC are based on OBR forecasts for GDP. They forecast an increase of 33.8\% of GDP per capita from 2022 to 2050. 



\subsection{Including uncertainty}


Let
\[
\hat{A} = \sum_{i=1}^N \lambda_i \Big[ \zeta \hat{X}_i + (\zeta-1) \hat{p}_i + v_i \Big] + \epsilon
\]
We then have
\[
\hat{A} = (\zeta-1) \frac{\sum_{i=1}^N C_i}{\text{GDP}} + \sum_{i=1}^N \lambda_i v_i + \epsilon
\]

